\documentclass[../main.tex]{subfiles}
\begin{document}

% Helper: labelled placeholder rectangle for figures not yet drawn.
\newcommand{\figplaceholder}[2][0.85\textwidth]{%
  \begin{tikzpicture}
    \draw[thick, dashed, fill=gray!8]
      (0,0) rectangle (#1,0.42\textwidth);
    \node[align=center, text width=\dimexpr#1-1cm\relax]
      at (\dimexpr#1/2\relax,0.21\textwidth)
      {\textbf{[FIGURE PLACEHOLDER]}\\[4pt]\itshape #2};
  \end{tikzpicture}%
}

% Helper: monospace verbatim box for prompt templates.
\lstdefinestyle{promptbox}{
  basicstyle=\ttfamily\footnotesize,
  backgroundcolor=\color{gray!6},
  frame=single,
  framesep=6pt,
  breaklines=true,
  columns=fullflexible,
  showstringspaces=false,
  keepspaces=true,
}


\chapter{Methodology}
\label{chap:methodology}

% =============================================================================
\section{Overview}
\label{sec:overview}
% =============================================================================

This chapter describes the proposed system for automatic Big Five (OCEAN)
personality recognition from written text. The central design problem is that a
language model prompted directly with a raw essay tends to predict personality
from \emph{surface style} --- register, sentence length, emotional intensity ---
rather than from stable psychological content. Surface style is a confound: an
expressive introvert and an expressive extravert produce superficially similar
text yet carry opposite labels.

The proposed system addresses this by interposing an explicit psychological
representation between the raw text and the classifier. The pipeline has three
components, illustrated in Figure~\ref{fig:pipeline}:

\begin{enumerate}
  \item \textbf{Psychological Profiler.} Each essay is mapped to a structured
    thirty-facet profile using the NEO-PI-R facet schema. This representation
    captures stable trait-level evidence rather than surface style.
  \item \textbf{Retrieval-Augmented Generation (RAG).} The profile is embedded
    and used to retrieve the most psychologically similar training exemplars from
    an offline-built index via a two-step hybrid search: approximate
    nearest-neighbour retrieval over profile embeddings followed by re-ranking
    with a trait-masked facet similarity score.
  \item \textbf{Reasoned Inference.} A language model receives the test essay
    and the retrieved exemplars --- rendered as full thirty-facet profiles grouped
    by trait --- and produces a \texttt{high}/\texttt{low} verdict through a
    structured reasoning chain.
\end{enumerate}

\begin{figure}[ht]
  \centering
  \figplaceholder[\textwidth]{%
    Figure 3.1 --- End-to-end pipeline. (1) Profiler: essay $\rightarrow$
    30-facet profile. (2) RAG: profile $\rightarrow$ embedding $\rightarrow$
    FAISS ANN candidate pool $\rightarrow$ re-rank by hybrid score
    $\rightarrow$ top-$k$ exemplars (all 30 facets + evidence + label).
    (3) Inference: test essay + full-profile exemplars $\rightarrow$ reasoned
    prompt $\rightarrow$ XML output $\rightarrow$ label.}
  \caption{The three-component pipeline. The profiler converts raw text into a
    structured psychological representation; the RAG component retrieves
    psychologically similar training exemplars via hybrid two-step search; the
    inference stage produces a reasoned \texttt{high}/\texttt{low} verdict per
    trait from full thirty-facet exemplar profiles.}
  \label{fig:pipeline}
\end{figure}

The pipeline operates in two phases. The \emph{offline phase} runs once over
the training corpus: every training essay is profiled, its profile is embedded,
and a similarity index is built. The \emph{online phase} runs once per test
essay: the essay is profiled label-blind, the index is queried, and the
retrieved exemplars are fed to the language model for classification.


% =============================================================================
\section{The Thirty-Facet Psychological Profiler}
\label{sec:profiler}
% =============================================================================

\subsection{Motivation}

A naive approach embeds the raw essay and retrieves nearest neighbours. Such a
retriever is dominated by lexical and stylistic similarity: it places two essays
close together when they share vocabulary, topic and register, with no notion
that \emph{warmth}, \emph{anxiety} or \emph{order} are the dimensions relevant
to personality. The profiler breaks this dependence by converting each essay
into a fixed psychological schema via an LLM-prompted facet extraction step
(GPT-4o); it is this schema, not the raw text, that drives both retrieval
and classification.

\subsection{The NEO-PI-R Facet Schema}
\label{subsec:neo-pi-r}

The schema is the facet structure of the Revised NEO Personality
Inventory~\cite{costa1992neo}. Each of the five OCEAN traits is decomposed into
six narrower facets, giving thirty facets in total (Table~\ref{tab:facets}).
Facets are more discriminative than a single trait score: two writers may both
be moderately extraverted, yet one is gregarious-but-passive and the other
assertive-but-solitary --- a distinction the six Extraversion facets capture but
a single score discards.

\begin{longtable}{p{1.4cm}p{3.0cm}p{2.6cm}p{6.0cm}}
  \caption{The thirty NEO-PI-R facets used by the profiler. Each facet receives
    one signal in $\{\textit{high},\textit{mod},\textit{low},\textit{none},
    \textit{n/e}\}$.}
  \label{tab:facets}\\
  \toprule
  \textbf{Code} & \textbf{Facet} & \textbf{Parent trait} & \textbf{Gloss} \\
  \midrule
  \endfirsthead
  \multicolumn{4}{c}{\textit{Table~\ref{tab:facets} continued}}\\
  \toprule
  \textbf{Code} & \textbf{Facet} & \textbf{Parent trait} & \textbf{Gloss} \\
  \midrule
  \endhead
  \midrule \multicolumn{4}{r}{\textit{continued on next page}}\\
  \endfoot
  \bottomrule
  \endlastfoot
  N1 & anxiety           & Neuroticism       & free-floating worry, fear, tension \\
  N2 & hostility         & Neuroticism       & anger, irritation, resentment toward others \\
  N3 & depression        & Neuroticism       & sadness, hopelessness, low mood \\
  N4 & self-conscious    & Neuroticism       & shame, embarrassment, social discomfort \\
  N5 & impulsiveness     & Neuroticism       & difficulty resisting urges or delaying \\
  N6 & vulnerability     & Neuroticism       & feeling overwhelmed, helpless under stress \\
  E1 & warmth            & Extraversion      & affectionate, friendly engagement with others \\
  E2 & gregariousness    & Extraversion      & preference for company, group settings \\
  E3 & assertiveness     & Extraversion      & leadership, dominance, decisive voice \\
  E4 & activity          & Extraversion      & energetic pace, busy schedule, motion \\
  E5 & excitement-seek   & Extraversion      & thrill-seeking, novelty appetite \\
  E6 & positive emotion  & Extraversion      & joy, enthusiasm, optimism \\
  O1 & fantasy           & Openness          & vivid imagination, daydreaming \\
  O2 & aesthetics        & Openness          & appreciation of art, beauty, nature \\
  O3 & feelings          & Openness          & emotional depth, introspective awareness \\
  O4 & actions/variety   & Openness          & openness to new experiences, change \\
  O5 & ideas             & Openness          & intellectual curiosity, abstract thought \\
  O6 & values            & Openness          & willingness to re-examine values, norms \\
  A1 & trust             & Agreeableness     & belief in others' good intentions \\
  A2 & straightforward   & Agreeableness     & candor, frankness, low manipulation \\
  A3 & altruism          & Agreeableness     & active concern for others' welfare \\
  A4 & compliance        & Agreeableness     & deference, low antagonism in conflict \\
  A5 & modesty           & Agreeableness     & humility, low self-promotion \\
  A6 & tender-minded     & Agreeableness     & sympathy, empathy with others' pain \\
  C1 & competence        & Conscientiousness & sense of capability, preparedness \\
  C2 & order             & Conscientiousness & tidiness, organization, structure \\
  C3 & dutifulness       & Conscientiousness & adherence to obligations, ethics \\
  C4 & achievement       & Conscientiousness & striving for goals, ambition \\
  C5 & self-discipline   & Conscientiousness & follow-through, resisting distraction \\
  C6 & deliberation      & Conscientiousness & careful planning, thinking before acting \\
\end{longtable}

\subsection{Signal Vocabulary}
\label{subsec:signal-vocab}

For every facet the profiler emits exactly one signal from a five-value
vocabulary: \texttt{high}, \texttt{mod}, \texttt{low}, \texttt{none}, and
\texttt{n/e} (``not evidenced''). The first three are ordinal levels of
evidence. \texttt{none} means the facet is observable in this genre but the
writer expresses it neutrally. \texttt{n/e} means the facet is structurally
unobservable in solo introspective writing --- interpersonal facets such as
\emph{warmth} or \emph{altruism} require social interaction the essay does not
record. Without a distinct not-evidenced value, the profiler would be forced to
fabricate signals on exactly those facets. This choice is consistent with
the linguistic-personality literature, which reports that interpersonal facets
of Agreeableness and Extraversion are the weakest to surface in written
self-report~\cite{pennebaker1999linguistic,mairesse2007using}.



% =============================================================================
\section{Profile Embedding and Trait-Aligned Retrieval}
\label{sec:retrieval}
% =============================================================================

\subsection{Profile Embedding}
\label{subsec:profile-embedding}

Once the profiler produces a structured profile, that profile is serialised to
text and passed through a sentence encoder to produce the retrieval vector. By
embedding the \emph{profile} rather than the raw essay, the retrieval space is
defined over the psychological facet vocabulary rather than surface lexis.

\subsection{Facet Vector Encoding}
\label{subsec:facet-encoding}

In parallel with the profile embedding, each profile is encoded as a point
in $\mathbb{R}^{30}$, one coordinate per facet. The categorical signal is
mapped to an ordinal value:
%
\begin{equation}
  \label{eq:signal-map}
  \texttt{high}\mapsto +1.0,\quad
  \texttt{mod}\mapsto +0.5,\quad
  \texttt{none}\mapsto 0.0,\quad
  \texttt{n/e}\mapsto 0.0,\quad
  \texttt{low}\mapsto -1.0.
\end{equation}
%
The scale is signed and symmetric; \texttt{mod} contributes only half the
weight of a confident signal; \texttt{none} and \texttt{n/e} both map to
zero so that absence of evidence neither attracts nor repels a candidate.
The resulting raw vector is L2-normalised:
%
\begin{equation}
  \label{eq:facet-vector}
  \boldsymbol{\phi}(p) \;=\; \frac{\mathbf{u}(p)}{\lVert \mathbf{u}(p) \rVert_2}.
\end{equation}

\subsection{Trait-Masked Facet Similarity}
\label{subsec:trait-slice}

When predicting a specific trait $\tau$, only the six facets of that trait are
relevant; the remaining twenty-four introduce cross-trait noise. Let
$S_\tau \subset \{1,\dots,30\}$ denote the six column indices belonging to
$\tau$. The trait-masked similarity between query $q$ and candidate $i$ is
%
\begin{equation}
  \label{eq:facet-cosine}
  \operatorname{sim}^{\mathrm{facet}}_{\tau}(q, i)
  \;=\;
  \cos\!\big(
    \boldsymbol{\phi}_q[S_\tau],\;
    \boldsymbol{\phi}_i[S_\tau]
  \big),
\end{equation}
%
where both sub-vectors are re-normalised after masking to remove any magnitude
bias from the discarded facets.

\subsection{Hybrid Retrieval Score}
\label{subsec:hybrid-score}

The retriever combines the dense profile-embedding similarity with the
trait-masked facet similarity into a single score. For query $q$, candidate
$d$, and trait $\tau$:
%
\begin{equation}
  \label{eq:hybrid}
  \operatorname{score}_\tau(q, d)
  \;=\;
  \beta_\tau \cdot \operatorname{sim}^{\mathrm{dense}}(q, d)
  \;+\;
  \gamma_\tau \cdot \operatorname{sim}^{\mathrm{facet}}_{\tau}(q, d).
\end{equation}
%
The weights $\beta_\tau$ and $\gamma_\tau$ are per-trait and satisfy
$\gamma_\tau > \beta_\tau$ for all traits, encoding the prior that the
structured facet signal is more trait-relevant than free-text similarity.
The gap is largest for Conscientiousness and Neuroticism --- the two traits
most susceptible to genre confound --- which therefore rely most heavily on
the structured channel.

Retrieval proceeds in two steps: an approximate nearest-neighbour search over
the profile-embedding index returns a candidate pool, which is then re-ranked
by the hybrid score to yield the final top-$k$ exemplars.

\subsection{Full-Profile Exemplar Rendering}
\label{subsec:exemplar-rendering}

Before being inserted into the prompt, each retrieved exemplar is rendered as
a \emph{full-profile}: all thirty facets are included, grouped by trait
(N\,/\,E\,/\,O\,/\,A\,/\,C), together with their evidence strings and a
linguistic summary block. The classifier therefore receives the complete
psychological portrait of each exemplar --- not raw essay text ---
alongside its known label.


% =============================================================================
\section{The Reasoned Inference Stage}
\label{sec:prompt}
% =============================================================================

\subsection{Structured Reasoning Format}

The language model is required to emit its output as six XML blocks in a fixed
order: \texttt{<evidence>}, \texttt{<facet\_check>}, \texttt{<cue\_tally>},
\texttt{<example\_alignment>}, \texttt{<verdict>}, and \texttt{<label>}. No
text is permitted outside these tags.

The \texttt{<verdict>} must be consistent with the \texttt{<cue\_tally>},
coupling the final answer to the audited evidence count rather than to the
tone of retrieved exemplars.

\subsection{Prompt Structure}

The user template provides the trait name, the HIGH and LOW definitions of the
trait, the retrieved full-profile exemplars (calibration anchors), and the test
essay. The template is shown in Listing~\ref{lst:reasoned}.

\begin{lstlisting}[style=promptbox,
  caption={Reasoned RAG user template.},
  label={lst:reasoned}]
Trait: {trait_name}

HIGH {trait_name}: {definition_high}
LOW  {trait_name}: {definition_low}
---

The following are the {top_k} most similar texts from the training set,
with their known labels and extracted psychological evidence. Use them
as calibration anchors - note how each one's cues map to its label.

{similar_context}

---

Text to classify:
<text>

Reason step by step and emit your output in the EXACT XML structure
specified in the system message.
\end{lstlisting}

Retrieval is performed per trait: a separate hybrid query is issued for each
of the five traits so that the masked facet score of
Equation~\eqref{eq:facet-cosine} is specialised to the trait being predicted.
Each retrieved exemplar is rendered as its full thirty-facet profile grouped by
trait block (N\,/\,E\,/\,O\,/\,A\,/\,C) with evidence strings, serving as a
calibration anchor for the classifier.
The default exemplar count is $k = 3$ per trait; the effect of $k$ is examined
in Chapter~\ref{chap:experiments}.


\end{document}
